% source: RH-Framework/theorems/T01_E4_NATIVE_FUNCTION_COMPLETION.md
\hypertarget{t01-e4-native-simple-function-tail-and-recognition-completions}{%
\chapter{T01-E4 --- Native Simple-Function Tail and Recognition Completions}\label{t01-e4-native-simple-function-tail-and-recognition-completions}}

\textbf{Classification:} \texttt{PROVED}

\textbf{Foundational origin:} \texttt{NATIVE\_DERIVED}

\textbf{Actual UGD function adapter:} \texttt{OPEN}

This theorem completes finite cut-scalar simple functions using the native
refinement content of T01-E1/E2/E3. It does not begin with an ordinary
\(L^1\), \(L^2\), Hilbert, Haar, or Lebesgue structure.

\begin{center}\rule{0.5\linewidth}{0.5pt}\end{center}

\hypertarget{finite-native-simple-functions}{%
\section{1. Finite native simple functions}\label{finite-native-simple-functions}}

Let \(R=(n_s,n_t,n_\alpha)\) be a native refinement grid and let
\(C_R\) be its cell atlas. A finite native simple function is a finitely
supported map

\[
f:C_R\to\mathbb K_\Sigma^0.
\tag{1.1}
\]

If two functions live on different grids, refine both to the coordinatewise
least-common-multiple grid. Refinement copies each coarse value to all fine cells
inside that coarse cell.

Addition, subtraction, dagger, and pointwise multiplication are defined after
common refinement.

\begin{center}\rule{0.5\linewidth}{0.5pt}\end{center}

\hypertarget{native-integral-gauges}{%
\section{2. Native integral gauges}\label{native-integral-gauges}}

The tail integral is

\[
\boxed{
\mathcal I_1(f)
=
\sum_{c\in C_R}\nu_R(c)D_\Sigma(f(c)).
}
\tag{2.1}
\]

The recognition integral is

\[
\boxed{
\mathcal I_2(f)
=
\sum_{c\in C_R}\nu_R(c)N_\Sigma(f(c)).
}
\tag{2.2}
\]

T01-E3 proves that both values are unchanged by refinement and integer cell
translation.

Define

\[
d_1(f,g)=\mathcal I_1(f-g),
\tag{2.3}
\]

\[
d_2(f,g)=\mathcal I_2(f-g).
\tag{2.4}
\]

The second expression is an energy gauge rather than a square-root norm.

\begin{center}\rule{0.5\linewidth}{0.5pt}\end{center}

\hypertarget{tail-gauge-control}{%
\section{3. Tail-gauge control}\label{tail-gauge-control}}

For cut scalars,

\[
D_\Sigma(x+y)
\le D_\Sigma(x)+D_\Sigma(y).
\]

Summing over a common refinement gives

\[
\boxed{
d_1(f,h)\le d_1(f,g)+d_1(g,h).}
\tag{3.1}
\]

Thus \(d_1\) is a rational refinement-invariant distance on finite simple
functions modulo equality after common refinement.

A sequence \((f_n)\) is \textbf{tail Cauchy} when for every positive cut rational
\(\varepsilon\),

\[
d_1(f_m,f_n)<\varepsilon
\]

for all sufficiently large \(m,n\).

Define

\[
\boxed{
\mathfrak F_{\Sigma,1}
=
\{\text{tail-Cauchy simple presentations}\}/\sim_{d_1}.
}
\tag{3.2}
\]

\hypertarget{theorem-3.1}{%
\subsection{Theorem 3.1}\label{theorem-3.1}}

\(\mathfrak F_{\Sigma,1}\) is a complete cut module. Refinement and integer cell
translation extend as isometries.

\hypertarget{proof}{%
\subsection{Proof}\label{proof}}

Addition and scalar multiplication are uniformly continuous by the scalar tail
triangle inequality and

\[
\mathcal I_1(zf)
=D_\Sigma(z)\mathcal I_1(f).
\]

The Cauchy-presentation quotient is well defined. Completeness follows from the
diagonal presentation construction using only positive cut-rational thresholds.
Refinement and translation preserve \(\mathcal I_1\) on the finite core, hence
extend to the completion. QED.

This object is the \textbf{native tail-integrable function completion}. Ordinary
\(L^1\) is not its definition.

\begin{center}\rule{0.5\linewidth}{0.5pt}\end{center}

\hypertarget{recognition-gauge-control}{%
\section{4. Recognition-gauge control}\label{recognition-gauge-control}}

For cut scalars,

\[
2N_\Sigma(x)+2N_\Sigma(y)-N_\Sigma(x+y)
=N_\Sigma(x-y)\ge0.
\tag{4.1}
\]

After summation,

\[
\boxed{
\mathcal I_2(f+g)
\le
2\mathcal I_2(f)+2\mathcal I_2(g).
}
\tag{4.2}
\]

A sequence \((f_n)\) is \textbf{recognition Cauchy} when for every positive cut
rational \(\varepsilon\),

\[
d_2(f_m,f_n)<\varepsilon
\]

for all sufficiently large \(m,n\).

Define

\[
\boxed{
\mathfrak F_{\Sigma,2}
=
\{\text{recognition-Cauchy simple presentations}\}/\sim_{d_2}.
}
\tag{4.3}
\]

\begin{center}\rule{0.5\linewidth}{0.5pt}\end{center}

\hypertarget{native-completed-pairing}{%
\section{5. Native completed pairing}\label{native-completed-pairing}}

For finite simple functions, define

\[
\boxed{
\langle f,g\rangle_{\Sigma,\nu}^0
=
\sum_c\nu_R(c)f(c)^\dagger g(c),
}
\tag{5.1}
\]

after common refinement.

The exact finite cut-square identity gives

\[
\boxed{
N_\Sigma\left(
\langle f,g\rangle_{\Sigma,\nu}^0
\right)
\le
\mathcal I_2(f)\mathcal I_2(g).
}
\tag{5.2}
\]

One proof clears the rational cell denominators and applies the finite
coefficient identity

\[
\left(\sum_iN(x_i)\right)
\left(\sum_iN(y_i)\right)
-
N\left(\sum_i x_i^\dagger y_i\right)
=
\sum_{i<j}N(x_iy_j-x_jy_i).
\tag{5.3}
\]

No pre-existing inner-product theorem is used.

\hypertarget{theorem-5.1}{%
\subsection{Theorem 5.1}\label{theorem-5.1}}

The pairing extends uniquely to

\[
\langle\cdot,\cdot\rangle_{\Sigma,\nu}:
\mathfrak F_{\Sigma,2}\times
\mathfrak F_{\Sigma,2}
\to
\widehat{\mathbb K}_\Sigma.
\tag{5.4}
\]

It satisfies dagger symmetry, cut-linearity, positivity, and

\[
\langle f,f\rangle_{\Sigma,\nu}=0
\Longleftrightarrow f=0.
\tag{5.5}
\]

\hypertarget{proof-1}{%
\subsection{Proof}\label{proof-1}}

Recognition-Cauchy presentations are energy bounded. Apply \texttt{(5.2)} to pairing
differences exactly as in T01-B. Positivity and definiteness follow from
\(\mathcal I_2(f)\) and the recognition-null equivalence relation. QED.

Therefore \(\mathfrak F_{\Sigma,2}\) is the complete \textbf{native recognition-
integrable function module}. Hilbert space may appear later only as a coordinate
shadow.

\begin{center}\rule{0.5\linewidth}{0.5pt}\end{center}

\hypertarget{native-bounded-simple-multipliers}{%
\section{6. Native bounded simple multipliers}\label{native-bounded-simple-multipliers}}

For a finite simple function \(h\), define its uniform cut bound

\[
\boxed{
U_\Sigma(h)=\max_c D_\Sigma(h(c)).
}
\tag{6.1}
\]

Then pointwise multiplication obeys

\[
\boxed{
\mathcal I_1(hf)
\le U_\Sigma(h)\mathcal I_1(f),
}
\tag{6.2}
\]

and

\[
\boxed{
\mathcal I_2(hf)
\le U_\Sigma(h)^2\mathcal I_2(f).
}
\tag{6.3}
\]

\hypertarget{proof-2}{%
\subsection{Proof}\label{proof-2}}

For every cell,

\[
D_\Sigma(h(c)f(c))
\le D_\Sigma(h(c))D_\Sigma(f(c))
\le U_\Sigma(h)D_\Sigma(f(c)),
\]

while

\[
N_\Sigma(h(c)f(c))
=N_\Sigma(h(c))N_\Sigma(f(c))
\le U_\Sigma(h)^2N_\Sigma(f(c)).
\]

Multiply by cell content and sum. QED.

Thus every bounded native simple multiplier extends continuously to both native
function completions.

This theorem does \textbf{not} claim that the product of two arbitrary tail-completed
functions remains tail-integrable. Such closure requires a separate hypothesis.

\begin{center}\rule{0.5\linewidth}{0.5pt}\end{center}

\hypertarget{explicit-cauchy-function}{%
\section{7. Explicit Cauchy function}\label{explicit-cauchy-function}}

On the unit-denominator grid, define

\[
f_N
=
\sum_{n=0}^{N}
2^{-(n+1)}1_{c_n},
\tag{7.1}
\]

where the \(c_n\) are distinct native cells. Then for \(M>N\),

\[
d_1(f_M,f_N)
=
\sum_{n=N+1}^{M}2^{-(n+1)}
<2^{-(N+1)},
\tag{7.2}
\]

and

\[
d_2(f_M,f_N)
=
\sum_{n=N+1}^{M}4^{-(n+1)}
<\frac{1}{3\,4^{N+1}}.
\tag{7.3}
\]

The same presentation therefore defines elements in both native completions.

\begin{center}\rule{0.5\linewidth}{0.5pt}\end{center}

\hypertarget{foundational-exclusion-boundary}{%
\section{8. Foundational exclusion boundary}\label{foundational-exclusion-boundary}}

The theorem does not use as definitions or proof generators:

\begin{verbatim}
ordinary L1 or L2 spaces;
Hilbert space;
Haar or Lebesgue integration;
operator norms;
complex floating-point values;
spectral or polar decomposition.
\end{verbatim}

The implementation uses exact cut scalars and rational cell content only.

\begin{center}\rule{0.5\linewidth}{0.5pt}\end{center}

\hypertarget{coordinate-shadow-boundary}{%
\section{9. Coordinate-shadow boundary}\label{coordinate-shadow-boundary}}

After the actual UGD refinement adapter is proved, one may attempt to identify

\begin{verbatim}
native tail function completion         with an ordinary L1 chart;
native recognition function completion with an ordinary L2 chart;
native bounded multipliers              with decomposable bounded operators.
\end{verbatim}

Those must be equivalence theorems. They may not be used as definitions of the
native objects above.

\begin{center}\rule{0.5\linewidth}{0.5pt}\end{center}

\hypertarget{exact-certificate}{%
\section{10. Exact certificate}\label{exact-certificate}}

The expected status is

\begin{verbatim}
PASS_NATIVE_SIMPLE_FUNCTION_TAIL_AND_RECOGNITION_COMPLETION_CORE
\end{verbatim}

The exact certificate verifies cross-refinement arithmetic, pairing identities,
native multiplier bounds, a nontrivial Cauchy function, and the absence of
classical numerical dependencies in the foundational module.

\begin{center}\rule{0.5\linewidth}{0.5pt}\end{center}

\hypertarget{remaining-gates}{%
\section{11. Remaining gates}\label{remaining-gates}}

\begin{verbatim}
T01-E4A NATIVE TAIL FUNCTION COMPLETION           PROVED
T01-E4B NATIVE RECOGNITION FUNCTION COMPLETION    PROVED
T01-E4C NATIVE BOUNDED SIMPLE MULTIPLIERS          PROVED

T01-E5 ACTUAL UGD REFINEMENT ADAPTER               OPEN
T01-E6 HAAR AND ORDINARY L1/L2 SHADOWS             OPEN
ACTUAL MP ACTIONS ON NATIVE FUNCTION DOMAINS       OPEN
NATIVE STATE-PACKET GENERATION                     OPEN
TARGET FAITHFULNESS                                OPEN
NATIVE SHUNYA                                      OPEN
RIEMANN HYPOTHESIS                                 OPEN
\end{verbatim}

\begin{flushright}\small\texttt{RH-Framework/theorems/T01\_E4\_NATIVE\_FUNCTION\_COMPLETION.md}\end{flushright}
